\section{Agentic Vision 的核心机制}

\subsection{从"静态观察"到"主动调查"}

传统 VLM 以\textbf{单一静态视角}处理视觉输入——如果遗漏了细微细节，只能被迫猜测。Agentic Vision 将图像理解从静态行为转变为\textbf{智能体过程}：

\begin{importantbox}{Agentic Vision 的三大创新}
\begin{itemize}
    \item \textbf{视觉推理 + 代码执行}：模型通过写 Python 代码来操作图像
    \item \textbf{多轮 Tool Call}：Think → Act → Observe 的 ReAct 循环
    \item \textbf{视觉草图（Visual Scratchpad）}：模型自主生成标注图来辅助推理
\end{itemize}
\end{importantbox}

\subsection{唯一的 Tool：Code Execution}

Agentic Vision 的 Tool 设计极其简洁——只有一个\textbf{代码执行工具}（Code Execution on Images）。通过写代码，模型可以实现：

\begin{itemize}
    \item \textbf{Bounding Box}：画检测框
    \item \textbf{Crop}：裁切图像区域
    \item \textbf{Rotate/Zoom}：旋转和放大
    \item \textbf{Annotate}：在图上打标注（文字、箭头、编号）
    \item \textbf{Draw}：绘制辅助线条
    \item \textbf{Save}：保存处理后的图像
\end{itemize}

\begin{warningbox}{关键洞察：BBox 不是外部 CV 模型产生的}
Agentic Vision 中的 Bounding Box 坐标是 \textbf{Gemini 2.5 Flash 自己产生的}——它在写代码时直接把 BBox 坐标传进去。这是模型\textbf{原生的基础视觉能力}（检测、分割、识别），而非依赖外部的目标检测模型。
\end{warningbox}

\subsection{ReAct 循环：Think → Act → Observe}

Agentic Vision 的工作流程是一个标准的 ReAct（Reasoning + Acting）过程：

\begin{enumerate}
    \item \textbf{Think}：分析用户问题，规划视觉操作步骤
    \item \textbf{Act}：生成 Python 代码，调用 Code Execution Tool
    \item \textbf{Observe}：查看代码执行结果（处理后的图像），决定是否需要进一步操作
    \item 循环直到得出最终答案
\end{enumerate}

\begin{knowledgebox}{这是 Post-Training 激励出来的}
ReAct 循环不是通过外部框架硬编码的，而是通过 Post-Training（强化学习）激励模型自发形成的多轮工具调用行为。模型学会了在什么时候该 zoom in、该 crop、该 annotate。
\end{knowledgebox}


